\documentclass[11pt]{article}
\usepackage[a4paper,margin=1in]{geometry}
\usepackage{amsmath,amssymb}
\usepackage[T1]{fontenc}
\usepackage{lmodern}
\usepackage{microtype}
\usepackage{xcolor}

\setlength{\parindent}{0pt}
\setlength{\parskip}{0.5em}

\newcommand{\promptbox}[1]{%
  \begingroup\setlength{\fboxsep}{10pt}\noindent
  \colorbox{blue!8}{\parbox{\dimexpr\linewidth-2\fboxsep\relax}{\textbf{Prompt. }#1}}
  \endgroup
}
\newcommand{\resultbox}[1]{%
  \begingroup\setlength{\fboxsep}{10pt}\noindent
  \colorbox{purple!8}{\parbox{\dimexpr\linewidth-2\fboxsep\relax}{\textbf{Result. }#1}}
  \endgroup
}
\newcommand{\examplebox}[1]{%
  \begingroup\setlength{\fboxsep}{10pt}\noindent
  \colorbox{green!8}{\parbox{\dimexpr\linewidth-2\fboxsep\relax}{\textbf{Example. }#1}}
  \endgroup
}

\title{\textbf{Book 1 --- Lecture 1 --- Interlude 1 --- Exercise 4}\\[2mm]
\large Magnitudes, Dot Product, and Angle}
\author{}
\date{}

\begin{document}
\maketitle

\promptbox{Let \(\mathbf{A}=(A_x,A_y,A_z)=(2,-3,1)\) and
\(\mathbf{B}=(B_x,B_y,B_z)=(-4,-3,2)\). Compute the magnitudes of
\(\mathbf{A}\) and \(\mathbf{B}\), their dot product, and the angle between
them.}

\section*{Computation}
\[
\begin{aligned}
\lVert\mathbf{A}\rVert &= \sqrt{2^2+(-3)^2+1^2}=\sqrt{14}\approx 3.741657,\\
\lVert\mathbf{B}\rVert &= \sqrt{(-4)^2+(-3)^2+2^2}=\sqrt{29}\approx 5.385165,\\
\mathbf{A}\cdot\mathbf{B} &= 2(-4)+(-3)(-3)+1(2)=3.
\end{aligned}
\]
Here \(\lVert\cdot\rVert\) denotes vector magnitude\footnote{Euclidean norm,
computed as square root of the sum of squared components.} and
\(\mathbf{A}\cdot\mathbf{B}\) denotes dot product\footnote{For
\(\mathbf{u},\mathbf{v}\in\mathbb{R}^n\), \(\mathbf{u}\cdot\mathbf{v}=\sum_i u_i
v_i\).}.

For the angle \(\theta\)\footnote{By convention, \(\theta\) often denotes an
angle between vectors in geometry and trigonometry.} between vectors:
\[
\cos\theta=\frac{\mathbf{A}\cdot\mathbf{B}}{\lVert\mathbf{A}\rVert\lVert\mathbf{B}\rVert}
=\frac{3}{\sqrt{14}\sqrt{29}}=\frac{3}{\sqrt{406}}.
\]
Hence,
\[
\theta=\arccos\!\left(\frac{3}{\sqrt{406}}\right)\approx 1.421388~\text{rad}
\approx 81.437539^\circ.
\]

\resultbox{\(\lVert\mathbf{A}\rVert=\sqrt{14}\), \(\lVert\mathbf{B}\rVert=\sqrt{29}\),
\(\mathbf{A}\cdot\mathbf{B}=3\), and the angle is
\(\theta\approx 81.44^\circ\).}

\examplebox{Since \(\mathbf{A}\cdot\mathbf{B}>0\), the angle is acute
(\(<90^\circ\)), matching the computed value \(81.44^\circ\).}

\end{document}
